\documentclass[aspectratio=169,11pt]{beamer}
\usepackage[spanish,es-noquoting,es-tabla]{babel}
\usepackage{fontspec}
\usepackage{amsmath}
\usepackage{booktabs}
\usepackage{tikz}
\usetikzlibrary{positioning,arrows.meta,calc}
\usepackage{pgfplots}
\pgfplotsset{compat=1.18}
% Misma paleta que el informe: las figuras compartidas la usan.
\definecolor{serie1}{HTML}{2A78D6}
\definecolor{serie2}{HTML}{EB6834}
\definecolor{serie3}{HTML}{1BAF7A}
\definecolor{serie4}{HTML}{EDA100}

\usetheme{default}
\usecolortheme[named=serie1]{structure}
\setbeamertemplate{navigation symbols}{}
\setbeamertemplate{footline}[frame number]
\setbeamerfont{frametitle}{series=\bfseries}
% Quién presenta cada bloque, en la esquina del título.
\newcommand{\presenta}[1]{\hfill{\normalfont\small\color{gray}#1}}

\title{Filtro RRCOS complejo de 8 coeficientes}
\subtitle{Implementaciones RTL en tiempo y en frecuencia}
\author{Andrés Villarreal \and Julián Moreyra \and Enzo Fernández}
\date{30 de octubre de 2026}

\begin{document}

\begin{frame}
  \titlepage
\end{frame}

% ---------------------------------------------------------------- Andrés
\begin{frame}{Qué se pidió \presenta{Andrés}}
  \begin{itemize}
    \item Filtro FIR complejo de 8 coeficientes, raíz de coseno elevado con \emph{roll-off} 0,5, sobre símbolos QPSK sobremuestreados 2×.
    \item Dos ramas: FIR en el dominio del tiempo y procesamiento por bloques en frecuencia (overlap-save con 50\,\% de superposición).
    \item En cada rama, una variante serie y una paralela, con la misma interfaz y los mismos resultados.
    \item Modelos en punto flotante y en punto fijo, comparados con distintos conjuntos de datos.
    \item Objetivos de reloj: 10\,MHz para las serie y 100\,MHz para las paralelas.
  \end{itemize}
\end{frame}

\begin{frame}{Cadena común a las dos ramas \presenta{Andrés}}
  \centering
    \begin{tikzpicture}[node distance=6mm and 8mm, >=Stealth,
      blk/.style={draw, rounded corners, minimum height=9mm, align=center, font=\small}]
    \node[blk] (gen) {Generador\\QPSK};
    \node[blk, right=12mm of gen] (up) {Sobremuestreo\\2×};
    \node[blk, right=12mm of up, yshift=8mm] (t) {FIR temporal};
    \node[blk, right=12mm of up, yshift=-8mm] (f) {Overlap-save\\FFT de 16 puntos};
    % Alto suficiente para recibir ambas ramas con flechas horizontales.
    \node[blk, minimum height=25mm, anchor=west] (cmp) at ($(f.east |- up)+(12mm,0)$) {Comparación\\y SQNR};
    \draw[->] (gen) -- node[above, font=\scriptsize]{$a[m]$} (up);
    \draw[->] (up.east) -- ++(4mm,0) |- (t.west);
    \draw[->] (up.east) -- ++(4mm,0) |- (f.west);
    \draw[->] (t.east) -- (t.east -| cmp.west);
    \draw[->] (f.east) -- (f.east -| cmp.west);
  \end{tikzpicture}


  \bigskip
  \begin{itemize}
    \item Los modelos en Python generan los coeficientes, la secuencia QPSK y las salidas esperadas de cada rama.
    \item En flotante las dos ramas difieren en menos de $2\times10^{-14}$: calculan el mismo filtro.
    \item El RTL recibe $x[n] \in \{-1, 0, +1\}$ en I y Q, con \texttt{valid}/\texttt{ready}/\texttt{last} y reset síncrono.
  \end{itemize}
\end{frame}

\begin{frame}{Punto fijo y SQNR \presenta{Andrés}}
  \begin{columns}[T]
    \begin{column}{0.55\textwidth}
      \resizebox{\linewidth}{!}{  \begin{tikzpicture}
    \begin{axis}[
        width=0.8\linewidth, height=6.5cm,
        xlabel={Bits fraccionarios del nodo variado}, ylabel={SQNR mínima (dB)},
        xmin=3.7, xmax=9.3, ymin=28, ymax=60, xtick={4,...,9},
        ymajorgrids, grid style={gray!20}, axis line style={gray!60}, tick style={gray!60},
        legend pos=north west, legend cell align=left, legend style={draw=none, font=\small},
        every axis plot/.append style={line width=1pt, mark size=2.2pt},
        /pgf/number format/use comma]
      \addplot[color=serie1, mark=*] coordinates {(5,29.8) (6,42.1) (7,44.3) (8,47.4) (9,57.6)};
      \addlegendentry{Tiempo: coeficientes y salida}
      \addplot[color=serie2, mark=square*] coordinates {(4,39.6) (5,43.0) (6,44.2) (7,44.5) (8,44.5) (9,44.5)};
      \addlegendentry{Frecuencia: \texttt{work}}
      \addplot[color=serie3, mark=triangle*] coordinates {(4,31.0) (5,33.6) (6,39.2) (7,41.4) (8,43.0) (9,42.9)};
      \addlegendentry{Frecuencia: $W$}
      \addplot[color=serie4, mark=diamond*] coordinates {(4,32.2) (5,37.4) (6,39.5) (7,41.0) (8,43.0) (9,42.3)};
      \addlegendentry{Frecuencia: $H$}
      \addplot[gray, dashed, line width=0.8pt, mark=none, forget plot] coordinates {(3.7,40) (9.3,40)};
      \node[gray, font=\small, anchor=north west] at (axis cs:6.6,39.8) {mínimo de 40\,dB};
    \end{axis}
  \end{tikzpicture}
}
    \end{column}
    \begin{column}{0.42\textwidth}
      \small
      \begin{itemize}
        \item Coeficientes Q1.7 y salida Q2.7 (9 bits) en las dos ramas, sin saturación: los anchos salen de cotas de peor caso.
        \item SQNR mínima con QPSK, semillas 0 a 4: \textbf{44,2\,dB} en tiempo y \textbf{43,0\,dB} en frecuencia, sobre un mínimo de 40\,dB.
        \item En tiempo el único error es cuantizar $h$; en frecuencia se suman $W$, $H$ y los redondeos internos.
      \end{itemize}
    \end{column}
  \end{columns}
\end{frame}

\begin{frame}{Verificación \presenta{Andrés}}
  \begin{itemize}
    \item Los modelos de punto fijo generan vectores con la entrada y la salida esperada de cada rama: impulso, ceros, extremos, bordes de trama y QPSK de 2048 muestras.
    \item Un banco de pruebas común verifica las cuatro variantes con los mismos vectores, en tres corridas: ideal, con pausas aleatorias en la entrada y la salida, y con reset en medio de una trama.
    \item La comparación es palabra por palabra y por índice de resultado, no por ciclo: vale para cualquier latencia.
    \item \texttt{make smoke} corre modelos, pruebas y simulación de cada variante registrada, igual en una máquina local y en GitHub Actions.
  \end{itemize}
\end{frame}

% ---------------------------------------------------------------- Julián
\begin{frame}{Tiempo serie \presenta{Julián}}
  \begin{itemize}
    \item Un sumador por componente recorre los 8 coeficientes: $y[n] = \sum_{k=0}^{7} h[k]\,x[n-k]$.
    \item Como $x \in \{-1, 0, +1\}$, cada producto es elegir $h_q[k]$, $-h_q[k]$ o 0: \textbf{sin multiplicador}, acumulación exacta en 9 bits.
    \item $II = 10$ ciclos por muestra (1 de aceptación, 8 de acumulación, 1 de entrega) y $L = 9$.
    \item 1\,Mmuestras/s a 10\,MHz.
  \end{itemize}
\end{frame}

\begin{frame}{Tiempo paralela: arreglo sistólico transpuesto \presenta{Julián}}
  \centering
  \begin{tikzpicture}[>=Stealth, node distance=4mm and 7mm, font=\small,
      reg/.style={draw, minimum width=8mm, minimum height=6mm},
      sel/.style={draw, rounded corners, minimum width=9mm, minimum height=5mm, font=\scriptsize},
      sum/.style={draw, circle, inner sep=1pt}]
    \node (x) {$x[n]$};
    \node[sel, right=8mm of x, yshift=10mm] (h7) {$\pm h_7$};
    \node[reg, below=4mm of h7] (s7) {$s_7$};
    \node[sum, right=8mm of s7] (a6) {$+$};
    \node[sel, above=4mm of a6] (h6) {$\pm h_6$};
    \node[reg, right=5mm of a6] (s6) {$s_6$};
    \node[right=5mm of s6] (dots) {$\cdots$};
    \node[sum, right=5mm of dots] (a1) {$+$};
    \node[sel, above=4mm of a1] (h1) {$\pm h_1$};
    \node[reg, right=5mm of a1] (s1) {$s_1$};
    \node[sum, right=8mm of s1] (a0) {$+$};
    \node[sel, above=4mm of a0] (h0) {$\pm h_0$};
    \node[reg, right=5mm of a0] (y) {$y$};
    \draw[->] (h7) -- (s7);
    \draw[->] (s7) -- (a6); \draw[->] (h6) -- (a6); \draw[->] (a6) -- (s6);
    \draw[->] (s6) -- (dots); \draw[->] (dots) -- (a1);
    \draw[->] (h1) -- (a1); \draw[->] (a1) -- (s1);
    \draw[->] (s1) -- (a0); \draw[->] (h0) -- (a0); \draw[->] (a0) -- (y);
    \draw[->] (y) -- ++(8mm,0) node[right] {$y[n]$};
    % x llega a todas las selecciones a la vez.
    \draw (x) |- ($(h7.north)+(0,4mm)$) -| (h0.north);
    \foreach \h in {h7,h6,h1} \draw[->] (\h.north |- {$(h7.north)+(0,4mm)$}) -- (\h.north);
    \draw[->] ($(h0.north)+(0,4mm)$) -- (h0.north);
  \end{tikzpicture}

  \bigskip
  \begin{itemize}
    \item Se guardan sumas parciales: $y[n] = h_0\,x[n] + s_1$, \; $s_k \leftarrow h_k\,x[n] + s_{k+1}$.
    \item Camino crítico: una selección y una suma de 9 bits, sin importar la cantidad de coeficientes.
    \item $II = 1$, $L = 1$; sin máquina de estados ni contador. 132 flip-flops (63 en la serie).
  \end{itemize}
\end{frame}

% ---------------------------------------------------------------- Enzo
\begin{frame}{Frecuencia serie: overlap-save con un multiplicador \presenta{Enzo}}
  \begin{columns}[T]
    \begin{column}{0.5\textwidth}
      \begin{itemize}
        \item Ventana de 16 muestras: 8 anteriores y 8 nuevas.
        \item FFT de 16 puntos, producto por $H[k]$, IFFT y salida de las posiciones 8 a 15.
        \item Un único multiplicador real compartido por FFT, producto espectral e IFFT.
        \item $L = 559$ ciclos; 0,139\,Mmuestras/s a 10\,MHz.
      \end{itemize}
    \end{column}
    \begin{column}{0.45\textwidth}
      \small
      \begin{tabular}{lr}
        \toprule
        Fase & Ciclos \\
        \midrule
        Copia y captura & 16 \\
        FFT (32 mariposas × 7) & 224 \\
        Producto espectral (16 × 6) & 96 \\
        Reordenamiento & 6 \\
        IFFT (32 mariposas × 7) & 224 \\
        Emisión y control & 9 \\
        \midrule
        Total por bloque de 8 & 575 \\
        \bottomrule
      \end{tabular}
    \end{column}
  \end{columns}
\end{frame}

\begin{frame}{Frecuencia paralela \presenta{Enzo}}
  \begin{itemize}
    \item FFT e IFFT en núcleos separados de 4 etapas con 8 mariposas cada una; 16 posiciones espectrales simultáneas, cada una con su $H[k]$ fijo.
    \item Cada producto multiplica siempre por la misma constante: se implementa con sumas y restas de desplazamientos (CSD), por ejemplo $181\,a = 2^8 a - 2^6 a - 2^4 a + 2^2 a + a$. Son 110 productos reales y ningún multiplicador.
    \item Pipeline: espectro listo en el ciclo 6, producto en el 9, IFFT en el 16, salida en el 17. $L = 26$ ciclos.
    \item Dos copias de cada almacenamiento solapan captura, cálculo y emisión: un bloque de 8 cada 8 ciclos, \textbf{una muestra por ciclo}.
  \end{itemize}
\end{frame}

\begin{frame}{Grado de paralelización \presenta{Enzo}}
  \begin{itemize}
    \item A igual reloj, el throughput sube $\times 10$ en tiempo y $\times 71{,}875$ en frecuencia: el trabajo de cada muestra pasa de hacerse en secuencia a hacerse en paralelo.
    \item En frecuencia, el producto espectral de un bloque tarda 2 de los 8 ciclos disponibles: con 4 posiciones en 4 tandas se llegaba al mismo throughput.
    \item Se desplegaron las 16 porque la consigna pide la máxima paralelización de la etapa FFT y para fijar cada $H[k]$ en su posición, a costa de área.
  \end{itemize}
\end{frame}

% ---------------------------------------------------------------- Andrés
\begin{frame}{Cómo medimos PPA \presenta{Andrés}}
  \begin{itemize}
    \item Flujo ASIC común: OpenLane~2 con SKY130 HD (130\,nm), de la síntesis al ruteo y la extracción de parásitos.
    \item Timing por análisis estático en 9 corners: proceso típico, lento y rápido, por tres extracciones de parásitos.
    \item Una variante cumple su reloj si setup y hold pasan en los 9 corners. Las restricciones eléctricas (transición, carga, fanout) se informan aparte.
    \item Potencia con OpenSTA y la actividad de la netlist final simulada con el benchmark QPSK de 2048 muestras.
    \item Las corridas en Vivado sobre Artix-7 quedan como referencia FPGA.
  \end{itemize}
\end{frame}

\begin{frame}{Resultados en SKY130 \presenta{Andrés}}
  \centering
  \small
  \begin{tabular}{lrrrr}
    \toprule
    & \multicolumn{2}{c}{Tiempo} & \multicolumn{2}{c}{Frecuencia} \\
    \cmidrule(lr){2-3}\cmidrule(lr){4-5}
    & Serie & Paralela & Serie & Paralela \\
    \midrule
    Reloj (MHz) & 10 & 100 & 10 & 100 \\
    Área de celdas (µm²) & 4748 & 8413 & 57\,365 & 541\,302 \\
    Setup y hold en 9 corners & cumple & cumple & cumple & \textcolor{serie2}{no cumple} \\
    Peor slack de setup (ns) & $+57{,}96$ & $+2{,}27$ & $+56{,}13$ & $-3{,}28$ \\
    Throughput (Mmuestras/s) & 1 & 99,95 & 0,139 & 98,75\textsuperscript{*} \\
    Potencia total (mW) & 0,087 & 2,23 & 0,658 & 63,5 \\
    Energía a 10\,MHz (nJ/muestra) & 0,087 & 0,022 & 4,73 & 0,643 \\
    \bottomrule
  \end{tabular}

  \medskip
  \footnotesize\textsuperscript{*}Condicionado: la frecuencia paralela no cumple 100\,MHz en los corners lentos (SS).
\end{frame}

\begin{frame}{Conclusiones \presenta{Andrés}}
  \begin{itemize}
    \item \textbf{Serie frente a paralela:} la paralela da una muestra por ciclo en las dos ramas y gasta menos energía por muestra (74\,\% menos en tiempo, 86\,\% en frecuencia), a costa de área.
    \item \textbf{Tiempo frente a frecuencia:} para este filtro gana la rama en tiempo en todo lo medido. Entre las paralelas, la de frecuencia ocupa 64 veces el área y consume 29 veces más energía por muestra.
    \item Con 8 coeficientes la FFT no se amortiza, y con entrada en $\{-1, 0, +1\}$ el FIR en tiempo no necesita multiplicadores.
    \item Pendiente: la frecuencia paralela no cierra 100\,MHz en los corners lentos. Una corrida preliminar con selector de salida one-hot reduce el peor setup a $-1{,}60$\,ns.
  \end{itemize}
\end{frame}

\begin{frame}{Planificación y aporte de cada uno \presenta{Andrés}}
  \centering
  \resizebox{!}{0.82\textheight}{% Días de octubre; x = (día - 3) * 0,45 cm. Fin exclusivo, como en docs/gantt.md.
  \begin{tikzpicture}[y=-0.42cm, font=\scriptsize]
    % Grilla semanal, fecha del registro, corte de C04 y entrega.
    \foreach \d in {5,12,19,26} {
      \draw[gray!25] ({(\d-3)*0.45},-1.3) -- ({(\d-3)*0.45},27.6);
      \node[gray, above] at ({(\d-3)*0.45},-1.3) {\d/10};
    }
    \draw[black!70] ({(21-3)*0.45},-1.3) -- ({(21-3)*0.45},27.6) node[below] {registro 20/10};
    \draw[black!70, dashed] ({(26-3)*0.45},-1.3) -- ({(26-3)*0.45},27.6) node[below] {corte C04};
    \draw[black!70, thick] ({(30-3)*0.45},-1.3) -- ({(30-3)*0.45},27.6) node[below] {entrega};
    % fila/tarea/plan inicio/plan fin/real inicio/real fin/color/terminada
    \foreach \y/\t/\ps/\pe/\rs/\re/\c/\ok in {
        0/E00/5/6/5/8/serie1/1, 1/E01/5/7/5/9/serie3/1, 2/E02/5/8/7/9/serie1/1,
        3/E03/8/11/9/10/serie2/1, 4/G01/5/30/3/21/serie1/0,
        6/A01/5/7/7/9/serie1/1, 7/A02/7/9/7/10/serie1/1, 8/A03/9/12/12/13/serie1/1,
        9/A04/9/13/12/13/serie1/1, 10/A05/13/15/12/14/serie1/1, 11/A06/15/20/14/21/serie1/1,
        12/A07/11/29/14/21/serie1/0,
        14/B01/8/11/12/14/serie2/1, 15/B02/11/16/14/15/serie2/1, 16/B03/11/14/15/17/serie2/1,
        17/B04/16/25/19/21/serie2/0,
        19/C01/7/12/6/14/serie3/1, 20/C02/12/20/12/16/serie3/1, 21/C03/12/14/15/19/serie3/1,
        22/C04/20/29/19/21/serie3/0,
        24/I01/29/30/0/0/serie1/0, 25/I02/14/30/15/21/serie1/0, 26/I03/20/30/0/0/serie1/0} {
      \node[anchor=east] at (-0.1,\y) {\t};
      \draw[gray!60, fill=gray!15] ({(\ps-3)*0.45},\y-0.42) rectangle ({(\pe-3)*0.45},\y-0.05);
      \ifnum\rs>0
        \ifnum\ok=1
          \fill[\c] ({(\rs-3)*0.45},\y+0.05) rectangle ({(\re-3)*0.45},\y+0.42);
        \else
          \fill[\c!45] ({(\rs-3)*0.45},\y+0.05) rectangle ({(\re-3)*0.45},\y+0.42);
        \fi
      \fi
    }
    % Grupos.
    \foreach \y/\g in {-0.1/Contrato y planificación, 5.9/Modelos y precisión (Andrés),
        13.9/Filtro temporal (Julián), 18.9/Filtro en frecuencia (Enzo), 23.9/Integración y presentación}
      \node[anchor=south west, font=\scriptsize\bfseries, inner sep=0] at (0,\y-0.55) {\g};
    % Leyenda.
    \begin{scope}[yshift=-29.2*0.42cm]
      \draw[gray!60, fill=gray!15] (0,0) rectangle (0.5,0.2); \node[right] at (0.5,0.1) {plan};
      \fill[serie1] (1.8,0) rectangle (2.3,0.2); \node[right] at (2.3,0.1) {Andrés};
      \fill[serie2] (3.8,0) rectangle (4.3,0.2); \node[right] at (4.3,0.1) {Julián};
      \fill[serie3] (5.6,0) rectangle (6.1,0.2); \node[right] at (6.1,0.1) {Enzo};
      \fill[gray!75] (7.2,0) rectangle (7.7,0.2); \node[right] at (7.7,0.1) {terminada};
      \fill[gray!30] (9.4,0) rectangle (9.9,0.2); \node[right] at (9.9,0.1) {en curso};
    \end{scope}
  \end{tikzpicture}
}
\end{frame}

\begin{frame}{Aprendizajes}
  \begin{itemize}
    \item \textbf{Andrés:} definir el contrato numérico antes del RTL ahorra discusiones después.
    \begin{itemize}\footnotesize
      \item Fijar formatos (Q1.7 / Q2.7), SQNR y vectores al principio (A04/A05) permitió verificar las cuatro variantes palabra por palabra con un único banco común.
      \item Una X en la simulación con retardos no siempre es una violación de timing: en Vivado aparecieron por los límites de pulso, las E/S sin restringir y un artefacto de CARRY4. Comparar contra la netlist sin SDF señal por señal llevó a la causa.
      \item Comparar PPA exige el mismo método para todos: plataforma, corners y criterio de «cumple». Pasar de Vivado a SKY130 a mitad de camino costó, pero dejó evidencia trazable.
    \end{itemize}
    \item \textbf{Julián:} \textcolor{gray}{pendiente}
    \item \textbf{Enzo:} \textcolor{gray}{pendiente}
  \end{itemize}
\end{frame}

\end{document}
